\section{Results}
\label{sec:results}

\subsection{Offline: joint decoding against cascades}
\label{sec:res-offline}

On the AMI reference windows (Table~\ref{tab:main7full}), \sys{} (offline) reaches
\num{v7:tx:en2zh:O2:cp} / \num{v7:tx:zh2en:O2:cp} cpBLEU$_\text{sess}$ (En$\rightarrow$Zh / Zh$\rightarrow$En) against
\num{v7:tx:en2zh:C-ft2:cp} / \num{v7:tx:zh2en:C-ft2:cp} for the cascade whose translator was trained on our windows
(exploratory comparison) and \num{v7:tx:en2zh:C-turn:cp} / \num{v7:tx:zh2en:C-turn:cp} for the one retrained on its own
diarizer turns (confirmatory: \sys{} ahead by \num{v7:reg:E-cas1:diff} and \num{v7:reg:E-cas2:diff}, $p_\text{Holm}$
\num{v7:reg:E-cas1:holm}). A speaker-blind model relabelled after decoding with the same diarizer reaches
\num{v6:relabel:en2zh:cp} / \num{v6:relabel:zh2en:cp}, close to \sys{} overall.

\subsection{Overlapped speech}
\label{sec:res-overlap}

\begin{table*}[t]
\centering\small
\setlength{\tabcolsep}{4pt}
\begin{tabular}{lrrrrrrrrr}
\toprule
 & \multicolumn{3}{c}{En$\rightarrow$Zh} & \multicolumn{3}{c}{Zh$\rightarrow$En} & \multicolumn{3}{c}{AliMeeting} \\
\cmidrule(lr){2-4}\cmidrule(lr){5-7}\cmidrule(lr){8-10}
System & all & overlap & CR & all & overlap & CR & all & overlap & CR \\
\midrule
\multicolumn{10}{l}{\emph{Reference windows, offline}} \\
Cascade, same data & 36.13 & 26.21 & 89.3 & 21.96 & 14.36 & 87.4 & 15.70 & 9.79 & 93.3 \\
Speaker-blind + relabel & 41.39 & 30.98 & 51.8 & 24.36 & 16.90 & 50.6 & 16.27 & 12.04 & 58.0 \\
\sys{} (offline) & 42.06 & 34.56 & 76.3 & 25.43 & 19.72 & 79.0 & 15.19 & 12.10 & 79.2 \\
\multicolumn{10}{l}{\emph{Whole meetings, online, 1\,s reads}} \\
Offline model + LA2 & 34.30 & 22.34 & 73.4 & 23.94 & 18.88 & 78.7 & 11.11 & 7.19 & 75.6 \\
Simul model + LA2 & 39.31 & 32.91 & 75.5 & 23.96 & 17.89 & 79.4 & 11.91 & 7.60 & 76.7 \\
Speaker-blind Simul + causal relabel & 35.22 & 27.17 & 65.5 & 19.33 & 9.06 & 54.4 & 12.67 & 8.55 & 63.0 \\
\sys{}-Simul & 40.09 & 35.71 & 82.3 & 24.53 & 17.73 & 80.2 & 11.47 & 9.18 & 80.2 \\
\bottomrule
\end{tabular}
\caption{Overlap-aware evaluation (\S\ref{sec:eval-overlap}): tcpBLEU over all units and over overlapped units, and
concurrent-speaker recall (CR, \%). Relabelling attaches one diarizer speaker to each translated segment. All rows and
metrics: Table~\ref{tab:overlap8}. AliMeeting references: machine translations of the human transcripts (Appendix~\ref{app:qe7}).}
\label{tab:ovmain8}
\end{table*}


Table~\ref{tab:ovmain8} splits the scores by overlap (\S\ref{sec:eval-overlap}). On single-speaker units relabelling
matches \sys{} (\num{v8:ov:en2zh:O2-nosd-relabel:1spk} vs.\ \num{v8:ov:en2zh:O2:1spk} in En$\rightarrow$Zh). On
overlapped units \sys{} leads by \num{v8:x:ov-O2-rel-en2zh:diff} (En$\rightarrow$Zh, 95\% interval
\num{v8:x:ov-O2-rel-en2zh:ci}) and \num{v8:x:ov-O2-rel-zh2en:diff} (Zh$\rightarrow$En, \num{v8:x:ov-O2-rel-zh2en:ci}),
both exploratory, and ties on AliMeeting (\num{v8:x:ov-O2-rel-ali:diff}); the gap opens with the amount of overlap
(Figure~\ref{fig:overlap8}). Relabelling gives one speaker per segment, so when two people talk at once one of them
disappears or is merged into the other (concurrent-speaker recall \num{v8:ov:en2zh:O2-nosd-relabel:cr}\%, collapse
\num{v8:ov:en2zh:O2-nosd-relabel:coll}\%); \sys{}, reading both speakers' activity, writes both (\num{v8:ov:en2zh:O2:cr}\%,
\num{v8:ov:en2zh:O2:coll}\%; recall gain \num{v8:x:cr-O2-rel-en2zh:diff} / \num{v8:x:cr-O2-rel-zh2en:diff} /
\num{v8:x:cr-O2-rel-ali:diff}, all $p<0.001$). Relabelling \sys{}'s own text shows that this gain is mainly
attribution, not better translation of overlapped speech (Appendix~\ref{app:v8}). The cascade, which translates every diarizer turn on its own, scores far
lower on overlapped units but recovers more concurrent pairs (\num{v8:ov:en2zh:C-ft2:cr}\%): \sys{} drops short
interjections (Appendix~\ref{app:v8}). With the human AliMeeting references, the per-direction streaming model ties a causally
relabelled speaker-blind model on overlapped units but has higher concurrent recall (\num{v9:ali:S:hum:cr} vs.\
\num{v9:ali:Rel:hum:cr}\%): on real Mandarin meetings reading the diarizer improves attribution of concurrent speakers,
not translation quality (Appendix~\ref{app:v9}).
\label{sec:res-real}
\label{sec:res-speech}

\subsection{Learning when to write and when to speak}
\label{sec:res-v10}

\begin{table*}[t]
\centering\small
\setlength{\tabcolsep}{4pt}
\resizebox{\textwidth}{!}{\begin{tabular}{llrrrrrrr}
\toprule
 & & \multicolumn{3}{c}{text tcpBLEU $\uparrow$} & speech & \multicolumn{3}{c}{latency (s) $\downarrow$} \\
\cmidrule(lr){3-5}\cmidrule(lr){7-9}
Test set & System & all & ov & 1spk & ASR-BLEU $\uparrow$ & text LAAL & speech LAAL & EndOffset \\
\midrule
\multirow{4}{*}{AMI Zh$\rightarrow$En} & Deadline SFT (strict baseline) & 25.02 & 19.97 & 27.10 & 18.07 & \textbf{2.99} & 4.01 & 4.08 \\
 & \quad + rl + PL-EM Talker & 25.02 & 19.97 & 27.10 & 20.35 & \textbf{2.99} & \textbf{3.61} & \textbf{3.13} \\
 & \sys{} v10 ($\bias=-1$) & \textbf{25.34} & \textbf{20.53} & \textbf{27.30} & \textbf{20.45} & 3.07 & 3.69 & 3.20 \\
 & SeamlessStreaming$^\dagger$ & 8.69 & 3.07 & 11.19 & 8.48 & 2.63 & 3.29 & 2.53 \\
\midrule
\multirow{4}{*}{AliMeeting Zh$\rightarrow$En} & Deadline SFT (strict baseline) & \textbf{21.43} & 16.91 & \textbf{23.05} & 16.73 & \textbf{2.82} & 4.28 & 4.57 \\
 & \quad + rl + PL-EM Talker & \textbf{21.43} & 16.91 & \textbf{23.05} & \textbf{18.08} & \textbf{2.82} & \textbf{3.65} & \textbf{3.23} \\
 & \sys{} v10 ($\bias=-1$) & 21.07 & \textbf{17.38} & 22.68 & 17.84 & 3.06 & 3.91 & 3.45 \\
 & SeamlessStreaming$^\dagger$ & 10.41 & 4.42 & 13.27 & 9.42 & 3.09 & 3.96 & 2.72 \\
\midrule
\multirow{4}{*}{AMI En$\rightarrow$Zh} & Deadline SFT (strict baseline) & 41.15 & 36.16 & 42.51 & 31.65 & \textbf{2.37} & 3.21 & 3.48 \\
 & \quad + rl + PL-EM Talker & 41.15 & 36.16 & 42.51 & 32.23 & \textbf{2.37} & \textbf{2.95} & \textbf{2.91} \\
 & \sys{} v10 ($\bias=-1$) & \textbf{41.74} & \textbf{36.48} & \textbf{43.24} & \textbf{32.34} & 2.47 & 3.05 & 2.99 \\
 & SeamlessStreaming$^\dagger$ & 13.28 & 5.86 & 16.71 & 11.31 & 2.31 & 3.19 & 3.04 \\
\bottomrule
\end{tabular}}
\caption{Test sets, each decoded once. Strict baseline: joint deadline-trained Thinker, frozen T1c Talker; second row: its text spoken with read-level merging (rl) and the PL-EM Talker; \sys{} v10: latent write alignment with $\bias=-1$ and the same speech layer. tcpBLEU over all / overlapped / single-speaker units; ASR-BLEU: mean of two TTS seeds; latency per speaker stream, chunk protocol (speech LAAL primary). Bold: best of our three systems. Tests: Table~\ref{tab:sig10}; $\bias=0$: Table~\ref{tab:delta0}. $^\dagger$External zero-shot single-speaker S2ST, no speaker attribution of its own: our windows, causal diarizer relabelling, one deterministic run (Appendix~\ref{app:e7}).}
\label{tab:main10}
\end{table*}


Table~\ref{tab:main10} gives the streaming systems on whole test meetings; all share training data, online windows and
diarizer. \emph{Speech layer.} Speaking the deadline baseline's own text with playback alignment (B+PA) raises ASR-BLEU
by \num{v10:sig:test:ami:speech:e2-b:diff} / \num{v10:sig:test:en:speech:e2-b:diff} /
\num{v10:sig:test:ali:speech:e2-b:diff} (AMI Zh$\rightarrow$En / En$\rightarrow$Zh / AliMeeting; all
$p_\text{Holm}<0.001$, Table~\ref{tab:sig10}) and shortens EndOffset by \num{v10:test:ami:e2-b:endoff} /
\num{v10:test:en:e2-b:endoff} / \num{v10:test:ali:e2-b:endoff}\,s, mainly by removing queueing (queue wait per unit
\num{v10:dev:z:b:p_queue_wait} vs.\ \num{v10:dev:z:e2:p_queue_wait}\,s, Zh$\rightarrow$En dev; Table~\ref{tab:decomp10}).
Speech starts hardly earlier (\num{v10:test:ami:e2-b:sstart} / \num{v10:test:en:e2-b:sstart} /
\num{v10:test:ali:e2-b:sstart}\,s): LAAL$_\text{spread}$ falls (\num{v10:test:ami:e2-b:slaal} /
\num{v10:test:en:e2-b:slaal} / \num{v10:test:ali:e2-b:slaal}\,s) because blocks wait less and play shorter. Under the conventional
block-start definition the dev gain of the full system shrinks to a few hundredths of a second
(Table~\ref{tab:m2laal10}).
\ifPostReview
On test, block-start speech LAAL is \num{v10:pr:test:ami:b:bslaal} / \num{v10:pr:test:ami:e2:bslaal} /
\num{v10:pr:test:ami:m:bslaal}\,s for B / B+PA / \sys{} on AMI Zh$\rightarrow$En (all sets, intervals, speech rate and
UTMOS: Appendix~\ref{app:postreview}).
\fi
B+PA differs from B in four ways at once (merging, the fine-tuned Talker, a backlog-dependent rate, best-of-two
CTC selection); on a fixed text, a rule-based layer with the same rate rule and selection gets most of the gain, and the
learned layer adds \num{v10:tk:ami:c-a:asr} / \num{v10:tk:ali:c-a:asr} ASR-BLEU and
\num{v10:tk:ami:c-a:slaal} / \num{v10:tk:ali:c-a:slaal}\,s of LAAL$_\text{spread}$ (AMI Zh$\rightarrow$En / AliMeeting,
Table~\ref{tab:talker10}). As selection uses a CTC recogniser and
scoring uses ASR, the gain may partly reflect ASR-friendly rather than more natural speech.

\emph{Write alignment.} On top of B+PA, the write-aligned Thinker (\sys{}) changes text tcpBLEU by
\num{v10:sig:test:ami:text:m-b:diff} / \num{v10:sig:test:en:text:m-b:diff} / \num{v10:sig:test:ali:text:m-b:diff}
(intervals \num{v10:sig:test:ami:text:m-b:ci} / \num{v10:sig:test:en:text:m-b:ci} / \num{v10:sig:test:ali:text:m-b:ci});
none is significant after Holm correction ($p_\text{Holm}$ \num{v10:sig:test:ami:text:m-b:holm} /
\num{v10:sig:test:en:text:m-b:holm} / \num{v10:sig:test:ali:text:m-b:holm}), and the AMI gains are of the size of the
seed spread of the per-direction deadline model (\num{v8:seed:zh2en:all} / \num{v8:seed:en2zh:all},
Appendix~\ref{app:v8}). It changes ASR-BLEU by only \num{v10:sig:test:ami:speech:m-e2:diff} /
\num{v10:sig:test:en:speech:m-e2:diff} / \num{v10:sig:test:ali:speech:m-e2:diff} (all intervals include zero) and is
heard later, by \num{v10:test:ami:m-e2:slaal} / \num{v10:test:en:m-e2:slaal} / \num{v10:test:ali:m-e2:slaal}\,s of
LAAL$_\text{spread}$, because it writes later (text LAAL \num{v10:test:ami:m-b:tlaal} /
\num{v10:test:en:m-b:tlaal} / \num{v10:test:ali:m-b:tlaal}\,s). In this study, write alignment thus yields no
significant improvement; the speech gains of \sys{} over B (\num{v10:sig:test:ami:speech:m-b:diff},
$p_\text{Holm}$ \num{v10:sig:test:ami:speech:m-b:holm}; \num{v10:sig:test:en:speech:m-b:diff}, $p_\text{Holm}$
\num{v10:sig:test:en:speech:m-b:holm}) are those of the playback layer. The AliMeeting text gap is a brevity penalty
(Appendix~\ref{app:dup}). As \sys{} writes later than B on every set, the AMI text differences may be a move along the
quality--latency curve; we did not run the controls that would tell (Limitations).

\paragraph{Ablation and the write bias.}
\label{sec:res-drift}
Without the bias, write alignment writes later than B+PA on the development set (text LAAL
\num{v10:dev:z:m0-e2:tlaal} / \num{v10:dev:e:m0-e2:tlaal}\,s; Table~\ref{tab:abl10}): as expected for a
likelihood-trained policy whose lattice can only defer, the posterior mass on the deadline plan falls over training from
\num{v10:drift:v2p:q_first} to \num{v10:drift:v2p:q_last} (Table~\ref{tab:drift10}). The bias undoes this; a
training-time latency budget acts less (Appendix~\ref{app:calib}).

\subsection{Test-set hygiene}
\label{sec:res-hygiene}

The final streaming numbers are not fully blind. (i) The write bias was introduced after a first test decoding without
it ($\bias=0$, italic rows of Table~\ref{tab:main10}), which wrote and spoke later than B, most on AliMeeting (text
tcpBLEU \num{v10:test:ali:m0:all} vs.\ \num{v10:test:ali:b:all}); its value was set on development meetings only, but
adding it was a test-informed decision that repaired what that run showed. (ii) The stage-2 checkpoint (step \num{v10:c:j:s2step}) is not the one the selection rule picked (step
\num{v10:c:j:s2rule}, never decoded on test): its dev5 criterion is higher (\num{v10:s2:3500:J} vs.\
\num{v10:s2:4000:J}) by a margin within decoder noise, and its En$\rightarrow$Zh text lag exceeds the rule's bound by
\num{v10:s2:e:3500:dL}\,s, a bound set relative to a later-writing system, not B (Appendix~\ref{app:v10train}). (iii) The same 16 AMI and 20
AliMeeting test meetings served throughout development, while the window closing gap, the latency protocol, read-level
merging, a short-turn fallback and the AliMeeting references changed; adaptive overfitting cannot be ruled out. (iv) The
analysis plans were internal, and most per-direction streaming tests were not run (Table~\ref{tab:reg9}).
